\begin{textsample}{POS Dim 1 – ma – Score 48.00 – ma/ma\_em\_38.txt}  \label{ex:f1_pos_064}
\textbf{Categoria} I : Tempo e Espaço Nesse espectro, Tempo e Espaço explicam os fenômenos nas ciências humanas porque permitem identificar contextos, sendo \textbf{categorias} difíceis de se dissociar.
O tempo é \textbf{um} desafio sobre o qual se debruçaram e se debruçam grandes pensadores de diversas áreas do conhecimento.
Na perspectiva filosófica, o tempo constitui -se \textbf{enquanto} estrutura da relação do \textbf{sujeito} com ele próprio e com o mundo.
Para Kant ( 2013 ), espaço e tempo são noções a priori \textbf{que} nos permitem compreender a experiência dos sentidos, utilizando espaço para identificar distâncias entre objetos reais, \textbf{enquanto} \textbf{que} tempo é utilizado para comparar intervalos entre eventos reais.
Na história, o tempo assume significados e importâncias variados.
O fundamental é compreender \textbf{que} não existe \textbf{uma} única noção de tempo, \textbf{que} ele não é homogêneo nem possui linearidade.
Nessa perspectiva, os estudantes precisam desenvolver noções de tempo \textbf{que} estão para além da dimensão cronológica e ganham diferentes sentidos, \textbf{que} vão do simbólico ao abstrato, destacando as noções de tempo em diferentes sociedades.
O conceito de tempo é entendido, \textbf{aqui}, como produção humana, no movimento dinâmico e histórico de \textbf{uma} concepção \textbf{que} parte da história de vida do \textbf{sujeito} quando ele \textbf{consegue} \textbf{assimilar} os processos simultâneos \textbf{que} ocorrem na duração de suas ações no decurso de sua existência.
No \textbf{que} se refere ao espaço, a compreensão deve contemplar suas dimensões histórica e cultural, transpondo suas representações \textbf{cartográficas} e, com isso, compreender o espaço associado aos arranjos dos objetos de diversas naturezas e, também, às movimentações de diferentes grupos, povos e sociedades.
Com isto, percebe -se \textbf{que} a espacialidade se dilata ou se comprime no tempo, conforme consideremos \textbf{um} período ou outro, nos quais se contraponham diferentes possibilidades de os \textbf{homens} se movimentarem no espaço.
Assim, \textbf{homem}, espaço e tempo aparecem como três fatores indissociáveis, não se limitando ao espaço ou ao território historiográfico \textbf{que} constitui \textbf{uma} região administrativa, geográfica ou de qualquer outro tipo, \textbf{implicando} a dissolução dos objetos historiográficos \textbf{que} não se ajustam a esses limites. \textbf{Categoria} II : Território e \textbf{Fronteira} Território é \textbf{uma} \textbf{categoria} das ciências humanas comumente associada à extensão de terra, a \textbf{uma} porção da superfície terrestre sob o domínio de \textbf{um} grupo, e suporte para nações, estados, países ( BNCC ).
Nesse território, há a consolidação de domínio de \textbf{fronteiras}, domínios políticos, administrativos, formação de unidades da federação, de cidades, de poder jurisdicional, de poder e soberania.
No contexto da globalização e território, há a transmissão de \textbf{um} sentimento de \textbf{que} o \textbf{homem} não possui território, pois, com as tecnologias da informação e comunicação, ele pode estar em vários locais ao mesmo tempo, ou seja, existe \textbf{uma} “ multiterritorialidade¹¹ ”, havendo \textbf{uma} sobreposição de territórios, em \textbf{um} processo constante de “ desterritorialização e reterritorialização”.¹² \textbf{Fronteira} é \textbf{uma} \textbf{categoria} das ciências humanas \textbf{que} expressa todo \textbf{um} contexto de soberania de \textbf{um} povo, de \textbf{uma} nação, de \textbf{uma} cultura, de costumes, de organização social e de saberes.
Milton Santos nos ensina \textbf{que} fronteira deriva de “ front, expressão militar \textbf{que} designa aquele espaço onde a guerra está sendo travada exatamente pelo controle do espaço.
O front transforma -se em \textbf{fronteira} e o espaço, em território ” ( GONÇALVES, 2004, p. 212 ).
Nesse cenário de globalização, as \textbf{fronteiras} são de grande relevância, são móveis, \textbf{perdem} a homogeneidade, se relativizam, tornando -se flexíveis, em virtude de acordos dos blocos econômicos, das \textbf{fronteiras} supranacionais e dos políticos e econômicos.
Partindo do contexto das \textbf{fronteiras} e traçando \textbf{um} paralelo entre várias sociedades presentes num mesmo território, observa -se o isolamento de comunidades estigmatizadas, como as áreas periféricas dos centros urbanos e metropolitanos, o \textbf{que}, de certa forma, setoriza a cidade, criando \textbf{fronteiras} de conhecimentos e saberes.
Cabe mencionar, ainda, \textbf{que} a globalização está cada dia mais desenvolvida, no sentido de \textbf{que} informações, produtos, tecnologias e outros elementos circulam cada vez mais.
No entanto, \textbf{fomenta} -se vertiginosamente a ideia de \textbf{que} ela é perversa, por aprofundar o abismo existente entre \textbf{fronteiras}.
Com essas \textbf{categorias}, verifica -se \textbf{que} existe \textbf{uma} concepção \textbf{teórica} da geografia crítica, a \textbf{dialética} \textbf{que} compreende os fenômenos naturais, sociais, políticos, econômicos, presentes em sociedades diversas, isto é, o espaço geográfico.
Assim, ensinar geografia é compreender o mundo \textbf{contemporâneo} e situar -se, posicionar -se como \textbf{sujeito} crítico, racional, buscando no estudante, o desenvolvimento integral de competências acadêmicas, críticas, éticas-valorativas, políticas e tecnológicas.
Portanto, no ensino médio, ao estudarmos essas \textbf{categorias}, devemos possibilitar aos estudantes a criticidade, para \textbf{que} possam identificar os processos formadores na busca pela investigação do espaço local ao mundo, transformando o local e a sociedade em \textbf{que} \textbf{vivem}. \textbf{Categoria} III : Indivíduo, Natureza, Sociedade, Cultura e Ética Somos seres eminentemente sociais e construímos formas de convivência e sobrevivência \textbf{que} são sempre sociais.
O indivíduo \textbf{que} afirmamos ser só existe porque há o olhar do outro para nos reconhecer e aceitar.
O \textbf{homem} é animal socioemocional por natureza, e essa condição de seres sociais está expressa na necessidade do outro para \textbf{viver} na coletividade e nos \textbf{exigir} habilidades socioemocionais \textbf{que} são essenciais para a convivência pacífica entre os indivíduos – principalmente entender o outro, ouvir, adaptar -se às situações e aprender a ser parte de \textbf{um} grupo, permitindo, assim, \textbf{que} os educandos continuem alimentando suas reflexões e o exercício da empatia, tolerância e respeito, para construção de \textbf{uma} sociedade sustentável.
Assim, por mais \textbf{que} cada ser humano tenha sua \textbf{singularidade}, as pessoas são essencialmente capazes de se organizar para \textbf{uma} vida em comum e de se governar.
O ser humano, além de ter a necessidade de conviver com o outro, tem diferentes maneiras de dominar e interferir na natureza para garantir a sua existência.
Sua cultura se manifesta na forma como supre suas necessidades básicas de subsistência relacionadas com a natureza, a partir do alimentar -se, vestir -se, criar seus artefatos, seus códigos de expressar ideias, sentimentos e valores.
Isso ratifica a ideia do \textbf{homem} como \textbf{um} ser cultural, \textbf{uma} vez \textbf{que} é o único ser capaz de produzir cultura, ou seja, a cultura como tudo aquilo \textbf{que} não é produzido pela natureza, mas pela inteligência, pela habilidade, intencionalidade e inspiração do \textbf{homem}.
É importante destacar a relevância de compreender a vida social a partir da noção de cultura ; isso ampliou a concepção de \textbf{homem}, possibilitando entendimento do círculo social, do conhecimento do outro e da diversidade dos povos.
Dessa forma, negar o mundo cultural seria reduzir o \textbf{homem} à natureza \textbf{que} o contém, mas \textbf{que} também é transformada por ele em \textbf{um} dinamismo \textbf{que}, gradativamente, o eleva.
Sendo assim, sempre \textbf{que} \textbf{uma} cultura for abertamente percebida, algo novo será conquistado, e \textbf{um} determinado modo de vida poderá se \textbf{ver} refletido no espelho multifacetado das civilizações.
Portanto, é necessária \textbf{uma} formação humanística capaz de promover no jovem o exercício da autocrítica, com base nos valores das dimensões ética e cultural, \textbf{que} possibilite a ele tornar -se \textbf{um} indivíduo autônomo, dotado de responsabilidade social, \textbf{que} preze os valores humanos e \textbf{que} aprenda a \textbf{viver} junto nesta “ aldeia global ”.
É \textbf{preciso} começar por conhecer a si próprio, numa espécie de viagem interior guiada pelo conhecimento, meditação e exercício de autocrítica.
Dentro da composição das áreas \textbf{que} integram o conhecimento humano, a ética é \textbf{um} dos ramos mais debatidos \textbf{hoje} em quase todos os círculos da sociedade brasileira.
Nunca se fez tão importante delinearmos \textbf{uma} eficiente e promissora educação ética aos nossos jovens, tendo em vista \textbf{um} processo de ensino \textbf{que} seja integrador e \textbf{que} contemple o ser humano em sua formação como \textbf{um} todo.
O ensino da ética nas escolas passou por trajetórias \textbf{que} se confundem com o ensino de filosofia no ensino médio.
A Lei nº 9.394/96, em seu artigo 36, indicava \textbf{que} o estudante, ao término do ensino médio, pudesse “ demonstrar ” ter o “ domínio dos conhecimentos de sociologia e filosofia ”, atrelados à noção de “ exercício da cidadania ”.
Essa concepção não permitia enxergar o desenvolvimento dos meios pelos quais se promoveria a consecução deste objetivo.
Foi a própria falta de determinação positiva desta lei \textbf{que} obrigou a elaboração de \textbf{um} outro dispositivo legal : a Lei nº 11.684/08, \textbf{que} tornou filosofia e sociologia obrigatórias para o ensino médio.
No contexto \textbf{atual}, a BNCC reforça a importância da discussão sobre a ética no ensino médio, principalmente em virtude do aprofundamento da complexidade de questões sociais, culturais e individuais.
Nesse sentido, o ensino da ética tem muito a oferecer aos nossos \textbf{discentes}, pois são grandes as possibilidades e contribuições \textbf{que} este ensino pode trazer ao aprimoramento dos estudos escolares.
Entre estas contribuições e discussões, podemos citar temas como a \textbf{teoria} necropolítica e a visão de \textbf{um} mundo em \textbf{que} vários países foram originados dentro de \textbf{um} contexto de violência contra corpos racializados e escravizados a fim de controlá -los, manipulá -los e/ou discipliná -los, a imposição da morte, do direito de matar e da guerra como instrumentos de poder, a biossimbiose, os debates acerca da bioética, as discussões sobre o \textbf{que} é e quais os limites da vida, a ética ambientalista, a preservação e manutenção da diversidade da fauna e da flora, bem como o cuidado com recursos naturais básicos para a sobrevivência humana, como a água, por exemplo.
Faz -se, portanto, \textbf{imprescindível} o ensino da ética de forma mais profissional e incisiva nas escolas para a construção de gerações futuras mais conscientes de seu papel ativo, de seu protagonismo na sociedade e de sua responsabilidade sobre os impactos de seus atos sobre si mesmo e sobre a vida dos outros com os quais convivemos em sociedade. \textbf{Categoria} IV : Política e Trabalho Trabalho > O branco açúcar \textbf{que} adoçará meu café nesta manhã de Ipanema não foi produzido por mim nem surgiu dentro do açucareiro por milagre. \textbf{Vejo} -o puro e afável ao paladar como beijo de moça, água na pele, flor \textbf{que} se dissolve na boca.
Mas este açúcar não foi feito por mim.
Este açúcar veio da mercearia da esquina e tampouco fez o Oliveira, dono da mercearia. > > Ferreira Gullar A estrofe acima citada é do escritor e poeta maranhense Ferreira Gullar \textbf{que}, por meio da arte de rima, descreve a \textbf{categoria} trabalho, problematizando as questões importantes ligadas ao mundo laboral na sociedade capitalista.
Em nosso cotidiano, nos deparamos sempre com a \textbf{categoria} trabalho, \textbf{que} é \textbf{um} tema polêmico, sempre \textbf{remetendo} a várias questões, como desemprego, avanço tecnológico, trabalhador polivalente, servidão na modernidade, precarização do trabalho, direitos trabalhistas, lutas de classes, política de estado, desigualdades sociais, saúde do trabalhador, salários, entre outras.
Na perspectiva sociológica, a \textbf{categoria} trabalho existe para satisfazer as necessidades humanas, \textbf{que} vão das mais \textbf{simples}, como alimentação, até as mais complexas, como lazer, crença e fantasia.
Ao longo da história, essa \textbf{categoria} teve significados e \textbf{encaminhamentos} diferentes, sendo compreendida como tortura e dominação ou usada para classificar pessoas.
Na Idade Moderna, com o surgimento do capitalismo, o trabalho deixou de ser \textbf{visto} como \textbf{uma} atividade repugnante e transformou -se em algo capaz de trazer \textbf{dignidade} ao ser humano.
Sendo assim, a ideologia capitalista passa a instituir a orientação para o trabalho como realização individual e social.
Entretanto, o \textbf{que} vai se intensificar, a partir do \textbf{século} XVIII, é a degradação e exploração do trabalhador.
É assim também durante \textbf{século} XIX, com a intensificação da produção industrial na sociedade de consumo.
Na \textbf{contemporaneidade}, novos desafios se apresentam, como avanço tecnológico, inteligência artificial e robótica, causando o temor do desemprego e a precarização das relações trabalhistas, tornando -se \textbf{um} dos grandes problemas sociais em nível global.
Segundo Forrester ( 1997 ), o desempregado, \textbf{hoje}, não é mais \textbf{um} objeto de \textbf{uma} marginalização provisória, ocasional, \textbf{que} atinge apenas alguns setores.
Agora, ele está às voltas com \textbf{uma} implosão geral, com \textbf{um} fenômeno comparável a tempestades, ciclones e tornados, \textbf{que} não visam ninguém em particular, mas aos quais ninguém pode resistir.
Ele é objeto de \textbf{uma} lógica planetária \textbf{que} \textbf{supõe} a supressão daquilo \textbf{que} se chama trabalho, vale \textbf{dizer}, empregos.
De acordo com os Parâmetros Curriculares Nacionais do ensino médio ( 2000 ), a \textbf{categoria} trabalho – juntamente com cidadania e cultura – é \textbf{uma} das \textbf{categorias} fundamentais das ciências sociais presentes no ensino médio porque permite, inicialmente, \textbf{que} alguns paradigmas \textbf{teóricos} \textbf{metodológicos} da sociologia, antropologia, política, e também da economia, do direito e da \textbf{psicologia}, sejam identificados, analisados, construídos e apropriados pelo estudante, pelo cidadão \textbf{que} frequenta a escola.
No campo de estudo da sociologia, o trabalho é \textbf{um} dos objetos de estudos dos clássicos – Karl Marx, Émile Durkheim e Max Weber.
Cada \textbf{um} desses pensadores vai buscar compreender as questões \textbf{que} envolvem a relação do trabalho com os aspectos da vida social, construindo \textbf{teorias}, \textbf{métodos}, diferentes explicações.
Todos elegeram o trabalho como \textbf{um} dos objetos científicos de seus estudos.
Fica \textbf{explícita} a importância da reflexão e o legado desses pensadores na compreensão dessa \textbf{categoria} nos dias \textbf{atuais}, a exemplo de Karl Marx ( como valor ), Max Weber ( como racionalidade capitalista ) ou como elemento de interação do indivíduo na sociedade, em Émile Durkheim.
Nessa perspectiva, para a construção de \textbf{um} jovem crítico e atuante na sociedade, na qual também vai vivenciar o mundo do trabalho na \textbf{contemporaneidade}, é necessário \textbf{que} os estudantes possam compreender e analisar a diversidade de papéis dos múltiplos \textbf{sujeitos} e seus mecanismos de atuação, além de identificar os projetos políticos e econômicos em disputa nas diferentes sociedades.
Destarte, é importante \textbf{frisar} a importância do conhecimento das \textbf{categorias} citadas, \textbf{visto} \textbf{que} é \textbf{preciso}, cada vez mais, \textbf{fomentar} na escola as questões \textbf{que} afligem a sociedade e os jovens \textbf{que} estão em processo de formação.
Isto, para poder garantir \textbf{uma} educação de qualidade, \textbf{que} crie condições de capacitar o repertório curricular para \textbf{uma} ampliação da leitura de mundo e da compreensão da vida em sociedade, elementos fundamentais para \textbf{uma} sociedade \textbf{justa}, democrática e igualitária.
Portanto, as áreas de conhecimento \textbf{que} compõem a área de Ciências Humanas e Sociais Aplicadas devem ter o seu espaço no currículo para garantir \textbf{que} princípios e valores humanos estejam sempre vivos na memória da sociedade.
Política Podemos observar, facilmente, \textbf{que} o número de pessoas debatendo sobre política tem aumentado em virtude da \textbf{crise} política social mundial e, com isso, alguns discursos e conceitos, como esquerda, direita, fascismo, liberalismo, comunismo, se elevaram a objetos de discussão.
É importante compreender como os estudantes de nossas escolas entendem a política e como a \textbf{disciplina} da área pode interferir, ajudando a melhorar a qualidade do debate e o processo de construção da consciência política dos estudantes.
A discussão sobre o tema política e seu ensino nas escolas é de grande relevância por se tratar da formação dos indivíduos para o pleno exercício da cidadania.
A LDB aponta \textbf{um} pressuposto básico quando sugere a necessidade da nossa educação promover \textbf{um} efetivo processo de socialização junto à escolarização.
Sabemos \textbf{que} não nascemos com o conhecimento pleno das leis, dos nossos direitos e deveres e da forma como os \textbf{atores} sociais se organizam dentro do Estado.
Por isso, a escola deve garantir o pleno debate e o ensino dos conceitos relacionados à política para, assim, oportunizar \textbf{um} conhecimento mais \textbf{aprimorado} sobre o tema.
Nesse aspecto, por meio de experiências compartilhadas, a escola deve estimular o debate regrado, incentivar a troca de informações e a reflexão crítica sobre as características e o modo de funcionamento das instituições e do próprio sistema político, propiciando, ainda, o desenvolvimento do senso de justiça, do respeito à diversidade, da tolerância e da solidariedade, competências importantes para criar relações democráticas efetivas. \textbf{Hoje}, compreende -se \textbf{que} o espaço escolar deve estimular o exercício da cidadania por meio de interferências, contribuições e construção de \textbf{uma} proposta de educação política coerente e adequada aos nossos tempos.
Para \textbf{uma} proposta como essa, será \textbf{imperativo} determinar \textbf{um} currículo, tendo objetivos claros e adequados e conteúdos pertinentes, sem deixar de observar o respeito e a tolerância aos direitos humanos.
Entende -se \textbf{que} \textbf{uma} verdadeira educação em política deve ser pensada e incentivada para extrapolar e ir para além dos muros da escola, promovendo, assim, o exercício de \textbf{uma} democracia participativa, com preceitos \textbf{baseados} em \textbf{uma} cidadania ativa, o \textbf{que} levaria a \textbf{uma} transformação dos nossos \textbf{discentes} e à transformação social. -- - ¹¹ CHELOTTI, Marcelo Cervo.
Reterritorialização e Identidade Territorial.
Disponível em. ¹² THÉRY, Hervé.
Globalização, Desterritorialização e Reterritorialização.
Disponível em.

% matched lemmas: aprimorar, aqui, assimilar, ator, atual, basear, cartográfico, categoria, conseguir, contemporaneidade, contemporâneo, crise, dialético, dignidade, discente, disciplina, dizer, encaminhamento, enquanto, exigir, explícito, fomentar, frisar, fronteira, hoje, homem, imperativo, implicar, imprescindível, justo, metodológico, método, perder, preciso, psicologia, que, remeter, simples, singularidade, sujeito, supor, século, teoria, teórico, um, ver, viver
\end{textsample}
